\documentclass[aspectratio=169, 11pt]{beamer}
\usepackage{fontspec}
\setmainfont{TeX Gyre Termes}
\setsansfont{TeX Gyre Heros}
\setmonofont{Latin Modern Mono}
\usepackage[spanish,es-nodecimaldot]{babel}
\usepackage{amsmath, amssymb, amsfonts, booktabs, graphicx}
\usepackage{tikz}
\usetikzlibrary{shapes.geometric, arrows.meta, positioning, calc, fit, shadows}
\usetheme{Madrid}
\usecolortheme{beaver}
\definecolor{unalRojo}{RGB}{148, 36, 36}
\definecolor{verdeOK}{RGB}{34, 139, 34}
\definecolor{naranjaAlerta}{RGB}{230, 126, 34}
\definecolor{rojoAlerta}{RGB}{192, 57, 43}
\definecolor{azulReactor}{RGB}{41, 128, 185}
\definecolor{verdeEntrada}{RGB}{39, 174, 96}
\setbeamercolor{structure}{fg=unalRojo}
\setbeamercolor{block title}{bg=unalRojo!15, fg=unalRojo}
\setbeamercolor{block body}{bg=unalRojo!5}
\setbeamercolor{alerted text}{fg=rojoAlerta}
\newcommand{\BSF}{\textit{Hermetia illucens}}
\newcommand{\NDIR}{\text{NDIR}}
\newcommand{\MRV}{\text{MRV}}
\setbeamertemplate{navigation symbols}{}
\setbeamertemplate{footline}[frame number]

\begin{document}

% ---------------------------------------------------------------------------
% Slide 1: La pregunta
% ---------------------------------------------------------------------------
\begin{frame}{Limitaciones y riesgos de sistemas inteligentes}
    \begin{center}
        \begin{tikzpicture}
            \node[rectangle, rounded corners=5pt, draw=unalRojo, fill=unalRojo!5, 
                  text width=14cm, align=justify, inner sep=12pt, font=\small] at (0,0) {
                \textbf{Desarrolle un an\'{a}lisis cr\'{i}tico de las limitaciones y riesgos asociados al uso de sistemas inteligentes en tecnolog\'{i}as de bioconversi\'{o}n de residuos. Incluya escenarios: ruralidad colombiana, falta de suministro el\'{e}ctrico y ciberseguridad.}
            };
        \end{tikzpicture}
    \end{center}
\end{frame}

% ---------------------------------------------------------------------------
% Slide 2: De medici\'{o}n pasiva a infraestructura cr\'{i}tica
% ---------------------------------------------------------------------------
\begin{frame}{De medici\'{o}n pasiva a infraestructura de control activa}
    \begin{center}
        \begin{tikzpicture}[scale=0.85, transform shape,
            etapa/.style={rectangle, rounded corners=6pt, minimum width=3.2cm, minimum height=1.3cm, align=center, font=\small},
            flecha/.style={-{Stealth}, line width=2pt}]

            \node[etapa, draw=gray, fill=gray!12] (pas) at (-5.5,0) {Medici\'{o}n pasiva\\Sensor aislado};
            \node[etapa, draw=naranjaAlerta, fill=naranjaAlerta!12] (act) at (0,0) {Control activo\\Sensores + red + IA};
            \node[etapa, draw=rojoAlerta, fill=rojoAlerta!12] (con) at (5.5,0) {Consecuencia\\Fallo = ceguera operativa};

            \draw[flecha, naranjaAlerta] (pas.east) -- (act.west);
            \draw[flecha, rojoAlerta] (act.east) -- (con.west);
        \end{tikzpicture}
    \end{center}

    \vfill
    \begin{itemize}
        \item El proceso biol\'{o}gico contin\'{u}a independientemente del estado del sistema de monitoreo.
        \item Anaerobiosis, mortalidad larval y picos de metano pueden desarrollarse sin detecci\'{o}n.
        \item Precisamente en escenarios de mayor estr\'{e}s operativo, el monitoreo es m\'{a}s necesario.
    \end{itemize}
\end{frame}

% ---------------------------------------------------------------------------
% Slide 3: Ruralidad colombiana
% ---------------------------------------------------------------------------
\begin{frame}{Ruralidad colombiana: conectividad y capacidad local}
    \begin{columns}[T]
        \begin{column}{0.48\textwidth}
            \begin{alertblock}{Infraestructura real}
                \begin{itemize}
                    \item Internet intermitente o inexistente en veredas.
                    \item Sin t\'{e}cnicos locales para calibrar o reparar sensores.
                    \item Soporte t\'{e}cnico urbano demorado e ineficiente.
                \end{itemize}
            \end{alertblock}
        \end{column}
        \begin{column}{0.48\textwidth}
            \begin{alertblock}{Consecuencia operativa}
                \begin{itemize}
                    \item Monitoreo remoto se interrumpe.
                    \item Sistema inteligente se convierte en \textit{sistema opaco}.
                    \item Dependencia de nube inaccesible deja al operador sin informaci\'{o}n.
                \end{itemize}
            \end{alertblock}
        \end{column}
    \end{columns}
\end{frame}

% ---------------------------------------------------------------------------
% Slide 4: Suministro el\'{e}ctrico
% ---------------------------------------------------------------------------
\begin{frame}{Suministro el\'{e}ctrico: continuidad como requisito de supervivencia}
    \begin{columns}[T]
        \begin{column}{0.48\textwidth}
            \begin{block}{Demanda energ\'{e}tica}
                \begin{itemize}
                    \item Sensores \NDIR{}: 24 h.
                    \item Microcontroladores y ventilaci\'{o}n forzada.
                    \item Cortes de luz de horas o d\'{i}as en campo colombiano.
                \end{itemize}
            \end{block}
        \end{column}
        \begin{column}{0.48\textwidth}
            \begin{alertblock}{Efecto en cascada}
                \begin{itemize}
                    \item Sensores dejan de medir, datos se pierden.
                    \item Ventiladores se apagan.
                    \item En reactor denso: pocas horas = anaerobiosis + mortalidad masiva.
                \end{itemize}
            \end{alertblock}
        \end{column}
    \end{columns}

    \vfill
    \begin{block}{Mitigaci\'{o}n}
        \centering
        Paneles solares fotovoltaicos + bater\'{i}as de respaldo = contingencia, no lujo.
    \end{block}
\end{frame}

% ---------------------------------------------------------------------------
% Slide 5: Ciberseguridad
% ---------------------------------------------------------------------------
\begin{frame}{Ciberseguridad: integridad de la cadena de custodia}
    \begin{columns}[T]
        \begin{column}{0.48\textwidth}
            \begin{alertblock}{Amenazas}
                \begin{itemize}
                    \item Dispositivos IoT con IP expuesta.
                    \item Manipulaci\'{o}n de registros de emisiones.
                    \item Acceso interno no controlado.
                \end{itemize}
            \end{alertblock}
        \end{column}
        \begin{column}{0.48\textwidth}
            \begin{block}{Capas de protecci\'{o}n}
                \begin{itemize}
                    \item TLS 1.2+ (cifrado en tr\'{a}nsito).
                    \item Certificados X.509 (autenticaci\'{o}n de dispositivos).
                    \item VLANs y firewalls (segmentaci\'{o}n de red).
                    \item Firmware firmado criptogr\'{a}ficamente.
                \end{itemize}
            \end{block}
        \end{column}
    \end{columns}

    \vfill
    \begin{alertblock}{Impacto en MRV}
        \centering
        VCS y Gold Standard exigen integridad de la cadena de custodia. Manipulaci\'{o}n exitosa = violaci\'{o}n de certificaci\'{o}n con consecuencias legales.
    \end{alertblock}
\end{frame}

% ---------------------------------------------------------------------------
% Slide 6: Calidad de datos y costo
% ---------------------------------------------------------------------------
\begin{frame}{Calidad de datos y barrera de entrada econ\'{o}mica}
    \begin{columns}[T]
        \begin{column}{0.48\textwidth}
            \begin{alertblock}{Riesgo metrol\'{o}gico}
                \begin{itemize}
                    \item Sensores con deriva, polvo, humedad.
                    \item Red neuronal aprende patrones falsos con apariencia de certeza.
                    \item Operador deja de hacer inspecciones visuales.
                \end{itemize}
            \end{alertblock}
        \end{column}
        \begin{column}{0.48\textwidth}
            \begin{alertblock}{Riesgo econ\'{o}mico}
                \begin{itemize}
                    \item Inversi\'{o}n inicial + licencias + capacitaci\'{o}n.
                    \item En proyectos comunitarios: m\'{a}s caro que el beneficio.
                    \item Sistema inteligente como \textbf{barrera de entrada}, no inclusi\'{o}n.
                \end{itemize}
            \end{alertblock}
        \end{column}
    \end{columns}
\end{frame}

% ---------------------------------------------------------------------------
% Slide 7: Dise\~{n}o robusto
% ---------------------------------------------------------------------------
\begin{frame}{Dise\~{n}o robusto: mitigaciones pr\'{a}cticas}
    \begin{center}
\begin{tikzpicture}[
    scale=0.82, transform shape,
    sol/.style={
        rectangle,
        rounded corners=5pt,
        draw=verdeOK,
        fill=verdeOK!10,
        minimum width=2.8cm,
        minimum height=1.2cm,
        align=center,
        font=\small,
        inner xsep=8pt,
        inner ysep=6pt
    },
    flecha/.style={
        -{Stealth[scale=1.2]},
        line width=1.5pt,
        verdeOK,
        shorten >=6pt,
        shorten <=6pt
    },
    node distance=1.4cm
]

\node[sol] (sol) {Paneles solares\\+ bater\'{i}as};
\node[sol, right=of sol] (lor) {LoRa / LoRaWAN\\bajo consumo};
\node[sol, right=of lor] (loc) {Modo local\\almacenamiento offline};
\node[sol, right=of loc] (cal) {Calibraci\'{o}n\\peri\'{o}dica};

\draw[flecha] (sol) -- (lor);
\draw[flecha] (lor) -- (loc);
\draw[flecha] (loc) -- (cal);

\end{tikzpicture}
    \end{center}

    \vfill
    \begin{block}{Principio rector}
        \centering
        La tecnolog\'{i}a debe servir como apoyo al criterio del operador, no como reemplazo.
    \end{block}
\end{frame}

\end{document}
